\section{October 5, 2026}
% Based on pages 7--9 of 6-7-week.pdf.

We now turn from one-step methods to linear multistep methods.
Instead of constructing new RK stages on every step, a multistep
method reuses solution values and derivatives from earlier times.
Throughout this discussion the grid is uniform:
\[
 t_n=t_0+n\tau,\qquad \tau>0,\qquad f_n=f(t_n,x_n).
\]
The fixed coefficients below are for a constant step size.

\subsection{From one-step methods to methods with memory}

A one-step method computes \(x_{n+1}\) from \(x_n\) through
\(x_{n+1}=\Psi^{t_{n+1},t_n}(x_n)\). A \(k\)-step method uses
the history \(x_n,\ldots,x_{n+k-1}\) to determine \(x_{n+k}\).
It therefore requires \(k\) starting values,
\[
 x_0,x_1,\ldots,x_{k-1}.
\]
Only \(x_0\) is supplied by the IVP. The remaining values can be
computed with a suitable one-step method, such as an RK method.
Their accuracy will enter the multistep convergence estimate.

\begin{example}[Explicit midpoint, or leapfrog]
 Integrating the ODE over two steps gives
 \[
  x(t_{n+2})-x(t_n)=\int_{t_n}^{t_{n+2}}f(t,x(t))\,dt.
 \]
 The midpoint quadrature rule suggests
 \begin{equation}\label{eq:leapfrog-oct-05}
  \boxed{x_{n+2}-x_n=2\tau f_{n+1}.}
 \end{equation}
 This is an explicit two-step method. It retains both \(x_n\)
 and \(x_{n+1}\), even though only one derivative appears in the
 formula. Equivalently,
 \(x_{n+1}=x_{n-1}+2\tau f_n\).
\end{example}

\begin{example}[Simpson's multistep formula]
 Applying Simpson quadrature to the same integral gives
 \begin{equation}\label{eq:simpson-lmm-oct-05}
  \boxed{x_{n+2}-x_n
       =\frac\tau3\bigl(f_n+4f_{n+1}+f_{n+2}\bigr).}
 \end{equation}
 It is implicit because \(f_{n+2}=f(t_{n+2},x_{n+2})\)
 depends on the new value. It is still a two-step method:
 three time levels occur, but only two starting values are needed.
\end{example}

The multistep midpoint formula \eqref{eq:leapfrog-oct-05} is
different from the one-step explicit midpoint RK method and from
the implicit midpoint Gauss method. Their names refer to related
quadrature ideas, but their update rules and stability properties
are different.

\subsection{General form and characteristic polynomials}

\begin{definition}[Linear multistep method]
 A \(k\)-step LMM has the form
 \begin{equation}\label{eq:lmm-general-oct-05}
  \sum_{j=0}^k\alpha_jx_{n+j}
    =\tau\sum_{j=0}^k\beta_j f_{n+j},
  \qquad \alpha_k\neq0.
 \end{equation}
 The coefficients are independent of \(n\) and \(\tau\).
 The method is explicit if \(\beta_k=0\), and implicit otherwise.
 We may normalize \(\alpha_k=1\) by dividing both sides by it.
\end{definition}

The word ``linear'' refers to the linear combination of values and
derivatives. An implicit LMM still requires a nonlinear solve when
\(f\) is nonlinear. We assume that the representation has no
unnecessary common factor in the two coefficient polynomials.

Define the shift operator on sequences by
\[
 Ex_n=x_{n+1},\qquad E^jx_n=x_{n+j},\qquad Ef_n=f_{n+1}.
\]
The two characteristic polynomials are
\begin{equation}\label{eq:lmm-polynomials-oct-05}
 \rho(\zeta)=\sum_{j=0}^k\alpha_j\zeta^j,
 \qquad
 \sigma(\zeta)=\sum_{j=0}^k\beta_j\zeta^j.
\end{equation}
Then \eqref{eq:lmm-general-oct-05} can be written compactly as
\begin{equation}\label{eq:lmm-shift-oct-05}
 \rho(E)x_n=\tau\sigma(E)f_n.
\end{equation}
Here \(\zeta\) is a polynomial variable; later, \(z=\tau\lambda\)
will denote the scaled eigenvalue in a stability test.

For the two examples,
\[
 \begin{array}{c|c|c}
  \text{method}&\rho(\zeta)&\sigma(\zeta)\\ \hline
  \text{leapfrog}&\zeta^2-1&2\zeta\\[2pt]
  \text{Simpson}&\zeta^2-1&\tfrac13(\zeta^2+4\zeta+1)
 \end{array}.
\]
This shift-operator form is also a finite linear filtering relation
between the solution sequence and its derivative sequence.

\subsection{Consistency error and order}

Insert a smooth exact solution into the formula and define its
unnormalized residual by
\begin{equation}\label{eq:lmm-residual-oct-05}
 \mathcal L_\tau[x](t_n)
  =\sum_{j=0}^k\alpha_jx(t_n+j\tau)
       -\tau\sum_{j=0}^k\beta_jx'(t_n+j\tau).
\end{equation}
This is \(\rho(E)x(t_n)-\tau\sigma(E)x'(t_n)\), where the shift
now acts on exact samples. It is the quantity denoted by \(L\) in
the handwritten notes.

\begin{definition}[Consistency order of an LMM]
 An LMM has order at least \(p\) if
 \[
  \mathcal L_\tau[x](t)=O(\tau^{p+1})
 \]
 for every sufficiently smooth test function \(x\). Its order is
 exactly \(p\) if the next coefficient does not vanish identically.
\end{definition}

Some conventions divide the residual by \(\tau\), in which case
an order-\(p\) method has normalized residual \(O(\tau^p)\).
We keep the unnormalized convention used in class. A residual is
not itself a global error; propagation through the recurrence must
also be controlled.

For leapfrog, expand about the midpoint \(t_{n+1}\):
\begin{align*}
 x(t_{n+2})&=x+\tau x'+\frac{\tau^2}{2}x''
                         +\frac{\tau^3}{6}x'''+\cdots,\\
 x(t_n)&=x-\tau x'+\frac{\tau^2}{2}x''
                         -\frac{\tau^3}{6}x'''+\cdots.
\end{align*}
Subtracting and cancelling \(2\tau x'\) gives
\begin{equation}\label{eq:leapfrog-error-oct-05}
 \mathcal L_\tau[x](t_n)
   =\frac{\tau^3}{3}x'''(t_{n+1})+O(\tau^5).
\end{equation}
Thus leapfrog has order two. Similarly, Simpson quadrature gives
\begin{equation}\label{eq:simpson-error-oct-05}
 \mathcal L_\tau[x](t_n)
   =-\frac{\tau^5}{90}x^{(5)}(t_{n+1})+O(\tau^7),
\end{equation}
so Simpson's multistep formula has order four.

\subsection{Algebraic order conditions}

Expanding \eqref{eq:lmm-residual-oct-05} about \(t=t_n\) gives
\[
 \mathcal L_\tau[x](t)
   =\gamma_0x(t)+\gamma_1\tau x'(t)+\gamma_2\tau^2x''(t)+\cdots,
\]
where the scalar coefficients in this lecture are
\begin{align*}
 \gamma_0&=\sum_{j=0}^k\alpha_j,\\
 \gamma_m&=\frac1{m!}\sum_{j=0}^k\alpha_jj^m
       -\frac1{(m-1)!}\sum_{j=0}^k\beta_jj^{m-1},
       \qquad m\geq1.
\end{align*}
The notation \(\gamma_m\) distinguishes these scalar residual
coefficients from the global-error functions \(C_m(t)\).

\begin{lemma}[LMM order conditions]
 The method has order at least \(p\) if and only if
 \begin{equation}\label{eq:lmm-order-conditions-oct-05}
  \boxed{\sum_{j=0}^k\alpha_j=0,\qquad
  \sum_{j=0}^k\alpha_jj^m
    =m\sum_{j=0}^k\beta_jj^{m-1},\quad m=1,\ldots,p.}
 \end{equation}
 For \(m=1\), the last sum is simply \(\sum_j\beta_j\).
\end{lemma}

\begin{proof}
 The conditions are exactly \(\gamma_0=\cdots=\gamma_p=0\), which are
 sufficient by Taylor's formula. Necessity follows by testing
 successively on constant, linear, and higher-degree polynomials.
\end{proof}

The first two conditions have the particularly useful form
\begin{equation}\label{eq:lmm-consistency-oct-05}
 \boxed{\rho(1)=0,\qquad \rho'(1)=\sigma(1).}
\end{equation}
An LMM satisfying these conditions is called \emph{consistent};
it has order at least one. Equivalently, the general order conditions
can be encoded as
\[
 \rho(e^z)-z\sigma(e^z)=O(z^{p+1})\qquad(z\to0).
\]
This follows by inserting the exponential test function into the
linear residual operator.

Before normalization there are \(2k+2\) coefficients; an overall
nonzero scaling leaves the method unchanged, so there are
\(2k+1\) effective parameters. Comparing this number with the
\(p+1\) order conditions suggests the formal upper bound \(p\leq2k\)
discussed in class. To justify the bound, suppose a nontrivial formula
were exact on all polynomials of degree at most \(2k+1\).
Hermite interpolation at the \(k+1\) distinct nodes allows arbitrary
values of both a polynomial and its derivative at those nodes.
The residual could vanish for all those choices only if every
\(\alpha_j\) and \(\beta_j\) were zero, contradicting
\(\alpha_k\neq0\). This formal bound imposes no stability condition.
The next lecture shows explicitly why even a consistent method of
high order may fail to converge.
